\section{Introduction and statement of the result}\label{sec:introduction}

The Sherrington--Kirkpatrick (SK) model was introduced by Sherrington and
Kirkpatrick~\cite{SK1975}
as a mean-field Ising spin glass with independent Gaussian pair
interactions. Parisi's replica-symmetry-breaking description uses a
hierarchy of order parameters~\cite{Parisi1979}. Guerra's interpolation
established the Parisi upper bound for the pressure~\cite{Guerra2003},
and Talagrand proved equality~\cite{Talagrand2006}.

Farhi, Goldstone, and Gutmann introduced the quantum approximate
optimization algorithm (QAOA) as alternating evolution under an objective
Hamiltonian and a transverse-field mixer, starting from the uniform
superposition~\cite{FGG2014}. For this model, Farhi, Goldstone, Gutmann, and Zhou
developed an exact description of its performance at every fixed depth in
the infinite-size limit and proved concentration at fixed
parameters~\cite{FGGZ2022}. Basso, Farhi, Marwaha, Villalonga,
and Zhou related that limit to QAOA on large-degree trees and conjectured
that increasing depth eventually attains the Parisi optimum
\cite[Section~7, Equation~(7.2)]{BFMVZ2022}.
Boulebnane et al.~\cite{BKL2026} proved a fixed-depth spin--boson
representation and used it to obtain numerical evidence for convergence
towards the Parisi value. Their representation provides a useful
Hilbert-space language for the limit. The construction here must also
realize bounded classical computations and eliminate longitudinal seed
controls before selecting one ordinary word.

Sels and Morone~\cite{SM2026} studied QAOA on the SK model using a
semiclassical truncated-Wigner approximation and compared its performance
with quantum data. Their numerical results and heuristic analysis give
further evidence for convergence to the Parisi value and suggest a depth
scaling.

The order of limits distinguishes this question from the fixed-instance
optimality proved by Farhi, Goldstone, and Gutmann~\cite[Section~VI,
Equation~(10)]{FGG2014}:
for each fixed instance, increasing the depth approaches its optimum.
That result does not provide a finite depth that works as the number
of spins tends to infinity. We instead choose a complete finite word for
each desired accuracy and then let the number of spins tend to
infinity. The quantum construction simulates a finite classical
calculation, using temporary Gaussian seed controls and then
removing them.

There are two ways to supply the classical calculation. Auffinger,
Chen, and Zeng proved that the zero-temperature Parisi order parameter
has infinitely many support points~\cite[Theorem~1]{ACZ2020}; Chen
proved full support on $[0,1)$ and a smooth
density~\cite[Theorem~1.3]{Chen2026}.
Our primary route uses the companion
\emph{Full support of the zero-temperature Sherrington--Kirkpatrick
order parameter}~\cite{SKSupport2026}. Its scalar diffusion yields a
finite function of independent Gaussian variables and a coefficient
formula for its value. We realize that function directly in the
rooted Hilbert space and identify the coefficient formula with its
energy. The imported scalar facts are full support of the zero-temperature
order parameter, the diffusion's consistency and value identities,
and a finite Gaussian coefficient identity.
Section~\ref{sec:classical-optimization} states these inputs before
constructing their rooted realization.

We also give an independent positive-temperature route.
It follows the Parisi-diffusion approach of
Montanari~\cite{Montanari2019}, whose optimization theorem assumes
interval support of the positive-temperature Parisi measure, and the
broader framework of El Alaoui, Montanari, and Sellke~\cite{AMS2020}.
Lopatto's full-support theorem~\cite{Lopatto2026} supplies that
positive-temperature support property. The finite-iteration
Gaussian and edge-energy analysis of El Alaoui, Montanari, and
Sellke~\cite{EAMS2023} then gives a near-optimal descendant formula.
This second route has its own quantitative and finite-degree
statements and does not use the zero-temperature companion.

\subsection{Model and quantifiers}

Let $(J_{ij})_{1\le i<j\le n}$ be independent standard real Gaussian random
variables. For $\sigma\in\{-1,1\}^n$, set
\begin{equation}\label{eq:sk-model}
 h_{n,J}(\sigma)=\frac1{\sqrt n}\sum_{i<j}J_{ij}\sigma_i\sigma_j,
 \qquad
 C_{n,J}=\frac1{\sqrt n}\sum_{i<j}J_{ij}Z_iZ_j,
 \qquad B_n=\sum_{i=1}^nX_i.
\end{equation}
Here $X,Y,Z$ are the Pauli matrices. The positive SK optimum is
\begin{equation}\label{eq:parisi-constant}
 \Pstar=\lim_{n\to\infty}\frac1n\E\max_{\sigma}h_{n,J}(\sigma).
\end{equation}
The ground-state Parisi formula and its attainment are due to
Auffinger and Chen~\cite{AC2017}. We use the covariance convention
$\xi(t)=t^2/2$; its exact relation to the unordered-pair model is
stated with the scalar input in Section~\ref{sec:classical-optimization}.
Appendix~\ref{sec:positive-temperature-alternative} also obtains this
limit from the positive-temperature formula and a log-sum comparison.

For $p\ge1$ and $\gamma,\beta\in\R^p$, define
\begin{equation}\label{eq:qaoa-state}
 \ket{\psi_{n,p,J}(\gamma,\beta)}
 =e^{-\ii\beta_pB_n}e^{-\ii\gamma_pC_{n,J}}
 \cdots e^{-\ii\beta_1B_n}e^{-\ii\gamma_1C_{n,J}}
 \ket{+}^{\otimes n}.
\end{equation}
All parameters are deterministic. Write
\begin{equation}\label{eq:qaoa-value}
 v_p(\gamma,\beta)=\lim_{n\to\infty}\frac1n
 \E\bra{\psi_{n,p,J}(\gamma,\beta)}C_{n,J}
 \ket{\psi_{n,p,J}(\gamma,\beta)},
 \qquad Q_p=\sup_{\gamma,\beta\in\R^p}v_p(\gamma,\beta).
\end{equation}
The fixed-parameter limit exists~\cite{FGGZ2022,BFMVZ2022}.

\begin{theorem}\label{thm:main}
For the zero-field SK model~\eqref{eq:sk-model},
\begin{equation}\label{eq:main-result}
 \lim_{p\to\infty}Q_p=\Pstar.
\end{equation}
Equivalently, for every $\varepsilon>0$ there exist a finite integer $p$ and
deterministic real vectors $\gamma,\beta\in\R^p$, independent of $n$ and $J$,
such that $v_p(\gamma,\beta)\ge\Pstar-\varepsilon$.
\end{theorem}

This proves the conjectured eventual Parisi optimality in the
fixed-parameter thermodynamic formulation of~\cite{BFMVZ2022}.
Measurement in the computational basis immediately gives $Q_p\le\Pstar$.
Appending a zero-angle layer gives $Q_{p+1}\ge Q_p$. Thus the theorem reduces
to constructing a finite word at every prescribed accuracy. The final word
may contain inverse gates, zero angles, and very large angles, all of which
are allowed in~\eqref{eq:qaoa-state}. A finite word in costs and mixers can
always be written in that alternating form by merging consecutive gates of
the same kind and inserting identity gates.

The theorem is stated in expected energy. The argument gives no
quantitative finite-size estimate, depth bound, or efficient procedure
for choosing the angles.

\subsection{How the construction works}

The proof separates a classical choice of target from its quantum
implementation. The intermediate setting is the large-degree
regular tree: each vertex carries an independent Gaussian label,
and a finite rooted calculation only reads labels within a fixed
radius. Normalized sums of centered child calculations become
jointly Gaussian coordinates. Section~\ref{sec:foundations} gives a
Hilbert-space realization of these coordinates and a continuous
quadratic expression for the limiting energy.

The scalar companion provides a bounded odd function of finitely
many independent Gaussians whose value approaches $P_*$. In
Section~\ref{sec:classical-optimization}, each new Gaussian is the
normalized aggregate of a suitable calculation one level below.
An orthogonality induction makes the joint Gaussian law exact.
The one-particle part of the resulting function gives its rooted
energy, so the scalar value formula passes to this realization
without a separate finite-degree edge expansion.

We next approximate all transmitted functions by bounded smooth
ones. A finite label assigns a \emph{type} to each vertex. Splitting
the label space into sufficiently small types lets every
aggregation omit the receiver's type, the types of all its path
ancestors, and a small reserved set of helper types, at arbitrarily
small energy cost. These exclusions prevent a compiled operation
from feeding back into its own parameters. They do not assert
independence of calculations using the same descendant subtree.
Proposition~\ref{prop:classical-near-optimum} gives the resulting
bounded finite formula, with its output written as $\sin(2\theta)$
for a bounded angle calculation $\theta$.

Theorem~\ref{thm:classical-synthesis} implements these angles using
cost and mixer gates together with a longitudinal field formed from
the Gaussian labels. To implement a neighbor sum, commutators combine
interactions along a receiver--sender--helper--sender walk. When the
walk returns to the same sender, the helper spin, held in its initial
state, supplies the desired term. We subdivide the sender type and
retain only walks whose two sender visits lie in the same part. This
keeps every returning walk, while the different-sender contribution
vanishes as the parts become small. The proof also controls the effects inside
the lower calculations, which are called repeatedly within a fixed
surrounding circuit. Helpers have arbitrarily small total probability;
discarding their output is included in the energy estimate.

Theorem~\ref{thm:seed-simulation} then removes the longitudinal
field from the already selected finite circuit. An imperfect cost echo
creates a small component with a large coherent displacement. Selective
spin rotations change that component while canceling their leading
action on the rest of the state; a long cost echo magnifies the change
into a finite longitudinal field. Separating its direction from the
circuit's existing Gaussian coordinates gives a seed that can be reused
at every requested seed gate. A separate one-particle calculation shows
that this Gaussian identification also preserves the rooted energy.
The elementary ensemble-control ideas are related to those of Li and
Khaneja~\cite{LiKhaneja2005}; Section~\ref{sec:seed-simulation}
proves the estimates for one fixed circuit, including its changing
insertion directions and the singular echo.

After all these choices there is one finite ordinary cost--mixer
word. Only at this point do we invoke the fixed-word tree-to-SK
identity of Basso et al.~\cite{BFMVZ2022}. The exact error allocation
and the order of choices are given in Section~\ref{sec:conclusion},
after the three constructions have been established.

\subsection{A regular MaxCut consequence}\label{sec:regular-maxcut}

For a finite simple graph $G=(V,E)$ with $E\ne\varnothing$, put
\[
 H_{\mathrm{MC}}(G)=\frac12\sum_{\{u,v\}\in E}(I-Z_uZ_v),
 \qquad B_G=\sum_{v\in V}X_v.
\]
For $a,b\in\R^p$, let $\langle\cdot\rangle^G_{p,a,b}$ denote expectation
with the gate order and signs of~\eqref{eq:qaoa-state}, using cost
$H_{\mathrm{MC}}(G)$, mixer $B_G$, angles $a,b$, and initial state
$\ket{+}^{\otimes |V|}$.  The expected fraction of edges cut by a
computational-basis measurement is
\[
 q_{p,G}(a,b)=\frac{\langle H_{\mathrm{MC}}(G)\rangle^G_{p,a,b}}{|E|}.
\]
Write $\operatorname{OPT}(G)=\operatorname{MaxCut}(G)/|E|$ for the maximum
cut fraction.

\begin{corollary}[Regular MaxCut in successive limits]
\label{cor:regular-maxcut}
For every $\eta>0$ there are an integer $p\ge1$, deterministic normalized
angles $\gamma,\beta\in\R^p$, and an integer $D_0\ge2$ such that, for every
integer $D\ge D_0$ and every finite simple unweighted $(D+1)$-regular graph
$G$ of girth greater than $2p+1$,
\begin{equation}\label{eq:high-girth-maxcut}
 q_{p,G}(2\gamma/\sqrt D,\beta)
 \ge \frac12+\frac{\Pstar-\eta}{\sqrt D}.
\end{equation}
The word $p,\gamma,\beta$ is independent of $D$, the graph size, and $G$;
the displayed usual MaxCut cost angles, unlike the normalized angles
$\gamma$, are scaled with the degree.

Let $G_{n,d}$ be uniform among simple labelled $d$-regular graphs on $n$
vertices, along admissible $n$ with $nd$ even.  For the same word, put
\[
 \Delta_D(G)=\operatorname{OPT}(G)
             -q_{p,G}(2\gamma/\sqrt D,\beta).
\]
The word may be chosen so that
\begin{equation}\label{eq:random-regular-maxcut}
 \lim_{D\to\infty}
 \limsup_{\substack{n\to\infty\\ n(D+1)\ \mathrm{even}}}
 \Pr_G\!\left[\sqrt D\,\Delta_D(G_{n,D+1})>\eta\right]=0.
\end{equation}
Here the inner limit keeps $D,p,\gamma,\beta$ fixed, and $\Pr_G$ is only
over the graph: $q_{p,G}$ already averages quantum measurement.
\end{corollary}

The complete proof is given in Section~\ref{sec:regular-maxcut-proof},
after the fixed-word tree interface and the SK construction.

The optimality comparison concerns random regular graphs in these
successive limits, not all high-girth graphs or fixed-degree optima.
No joint size--degree limit is asserted. Here $p$ counts QAOA rounds,
not a degree-uniform two-qubit gate depth; the cut fraction above
averages quantum measurement.

\subsection{Reading paths and conventions}

Section~\ref{sec:foundations} develops the fixed-word rooted
calculus. Section~\ref{sec:classical-optimization} imports the scalar
theorems and constructs the bounded excluded formula.
Sections~\ref{sec:classical-synthesis} and~\ref{sec:seed-simulation}
implement it and remove the seeds. Section~\ref{sec:conclusion}
assembles the ordinary word and proves Theorem~\ref{thm:main}.
Section~\ref{sec:regular-maxcut-proof} then proves the regular-MaxCut
consequence from this theorem and the fixed-word tree interface.

For the positive-temperature reading path,
Appendix~\ref{sec:positive-temperature-alternative} replaces the
scalar input and its rooted realization. It derives the same
classical conclusion independently, using the shared bounded
approximation and exclusion lemmas of
Section~\ref{sec:classical-optimization}. The quantum constructions
and final assembly are the same for both paths.

Physical system size is $n$, tree degree is $D+1$, and ordinary
circuit depth is $p$. The scalar mesh count in
Section~\ref{sec:classical-optimization} is $N$; top-generation
occupation operators appearing in the operator estimates are
identified by their operator context. Mixer angles are denoted by
$\beta$, inverse temperature by $\beta_0$, and the scalar Parisi
coefficient $\gamma(t)$ is distinct from the vector of cost angles.
Inner products are antilinear in the first argument.

All constants may depend on a fixed finite formula or word.
Uniformity in any approximation parameter is stated when required.
Degree limits are taken with all words and their angles fixed;
further approximation limits concern the resulting tree values.
